\documentclass[journal]{IEEEtran}

% Compile with XeLaTeX.
\usepackage{fontspec}
\usepackage{xeCJK}
\usepackage{cite}
\usepackage{amsmath,amssymb,amsfonts}
\usepackage{algorithmic}
\usepackage{graphicx}
\usepackage{textcomp}
\usepackage{xcolor}
\usepackage{booktabs}
\usepackage{multirow}
\usepackage{array}
\usepackage{url}

\setmainfont{Times New Roman}
\setCJKmainfont{SimSun}
\setCJKsansfont{SimHei}
\setCJKmonofont{FangSong}

\hyphenation{op-tical net-works semi-conduc-tor}

\begin{document}

\title{VK-RMD:一种XXXX}

\author{Eagle%
\thanks{稿件接收日期占位；修改日期占位。}
\thanks{作者单位与地址占位。}
\thanks{通信作者：姓名占位。}}

\markboth{IEEE Internet of Things Journal,~Vol.~xx, No.~xx, Month~2026}%
{作者等：中文标题占位}

\maketitle

\begin{abstract}
多模态人类感知技术已经在
\end{abstract}

\begin{IEEEkeywords}
关键词占位，关键词占位，关键词占位。
\end{IEEEkeywords}

\section{引言}
本文引言占位。

% 第一张总体框架图建议放在引言末尾：
% 先完成问题动机和贡献概述，再给出总体框架图。
\begin{figure*}[!t]
    \centering
    \fbox{\parbox[c][2.2in][c]{0.92\textwidth}{\centering 图 1 总体框架占位}}
    \caption{总体框架占位。}
    \label{fig:framework}
\end{figure*}

\section{相关工作}

\subsection{小节占位}

\section{系统模型与问题定义}

\subsection{问题定义}

\subsection{符号说明}

\section{方法}

\subsection{整体架构}
为实现低视觉暴露条件下的缺失模态人体感知，本文希望设计一个能够动态处理不同数量传感器输入的基础模型。与传统固定输入模型不同，真实部署中可用模态集合并不稳定：视觉关键点、深度、LiDAR、毫米波雷达和 WiFi-CSI 可能任意组合出现。因此，模型首先需要支持可变长度的模态输入；其次，不同传感器具有不同数据形态和噪声特性，模型需要保留每个模态的专属信息；最后，在多模态融合时，模型还应根据当前可用模态组合动态分配决策贡献，避免简单拼接或平均融合造成信息压缩和弱模态干扰。

受 X-Fi 的模态不变多模态人体感知思想启发，本文同样采用“模态专属编码器 + 跨模态 Transformer + 动态融合”的整体结构，但问题设定和训练目标不同。X-Fi 旨在学习一个可接收任意模态组合的模态不变人体感知基础模型，其核心是通过跨模态 Transformer 和 X-fusion 在不同传感器组合间学习共享表征。本文的目标不是重新提出一个通用模态不变 foundation model，而是在低视觉暴露约束下研究缺失模态鲁棒性：下游模型不使用原始 RGB，而以视觉关键点作为结构输入，并通过 Full-to-Subset Guidance 让完整模态教师约束随机缺失模态学生。

具体而言，本文框架包含三类模块。第一，模态专属编码器分别处理不同传感器输入。视觉关键点分支不再使用图像卷积网络，而是将关节坐标视为人体结构图进行编码；深度、LiDAR、毫米波雷达和 WiFi-CSI 则通过各自的传感器编码器提取模态专属特征。第二，所有模态特征被投影到统一的 body-token 空间，并加入模态嵌入，使模型既能在统一维度上处理不同模态，又能保留每个模态的来源信息。第三，融合模块根据当前可用模态子集生成模态贡献权重，并输出 HPE 姿态估计或 HAR 动作识别结果。

与 X-Fi 的关键区别在于训练关系。X-Fi 主要强调同一模型对不同模态组合的适配能力，而本文显式构造 teacher/student 输入条件不对称：teacher 使用完整模态集合学习完整感知行为，student 在随机采样的可用模态子集上训练，并学习逼近 teacher 的输出、结构 token 和融合行为。因此，本文的核心不是简单的模态 dropout，也不是传统大模型压缩小模型的知识蒸馏，而是 full-modality-to-subset guidance，即完整模态感知行为向缺失模态子集感知行为的迁移。

\subsection{模块占位}

\subsection{训练目标}

\begin{equation}
    \mathcal{L} = \mathcal{L}_{\mathrm{task}} + \lambda \mathcal{L}_{\mathrm{aux}} .
    \label{eq:objective}
\end{equation}

\section{实验设置}

\subsection{数据集与任务}

\subsection{基线方法}

\subsection{评价指标}

\subsection{实现细节}

\section{结果与分析}

\subsection{主结果}

\begin{table*}[!t]
    \centering
    \caption{主结果表占位。}
    \label{tab:main_results}
    \begin{tabular}{lcccc}
        \toprule
        方法 & 设置 & 指标 1 & 指标 2 & 指标 3 \\
        \midrule
        占位 & 占位 & -- & -- & -- \\
        \bottomrule
    \end{tabular}
\end{table*}

\subsection{消融实验}

\begin{table}[!t]
    \centering
    \caption{消融实验表占位。}
    \label{tab:ablation}
    \begin{tabular}{lcc}
        \toprule
        变体 & 指标 1 & 指标 2 \\
        \midrule
        占位 & -- & -- \\
        \bottomrule
    \end{tabular}
\end{table}

\subsection{诊断分析}

\section{讨论}

\section{结论}

\section*{致谢}
致谢占位。

\section*{数据与代码可用性}
数据与代码可用性占位。

\begin{thebibliography}{1}

\bibitem{ref_placeholder}
参考文献占位。

\end{thebibliography}

\end{document}
